The supplied solution formulates the perpendicular-intercept geometry
as a minimum-time interception. Figure 2 shows the absolute rocket path
$NC$ perpendicular to the asteroid path $AC$.
\ioaafig{2022_TH_Q12-s1.pdf}
With $L=NA$, the distances traveled are
\[
D=v_At_{\min}=L\cos\alpha,\qquad d=v_Rt_{\min}=L\sin\alpha.
\]
Thus
\[
t_{\min}=\frac{L\cos\alpha}{v_A}=\frac{L\sin\alpha}{v_R},
\qquad v_R=v_A\tan\alpha.
\]
The source then prints
\[
\vec v_{RN}=\vec v_R-\vec v_A,
\qquad
v_{RN}^2=v_R^2+v_A^2-2v_Rv_N\cos(90^\circ-\alpha)
=v_R^2+v_A^2-2v_Rv_N\sin\alpha.
\]
% sic: The printed relative-velocity subtraction and squared term use v_A where v_N would be expected for the spacecraft velocity.
The transverse components in the figure give
\[
v_R\cos\alpha=v_{RN}\cos(\beta-90^\circ)=v_{RN}\sin\beta,
\]
so
\[
\sin\beta=\frac{v_R}{v_{RN}}\cos\alpha
  =\frac{v_A}{v_{RN}}\sin\alpha.
\]
Substituting $v_R=v_A\tan\alpha$ into the source's squared-speed expression,
\begin{align*}
v_{RN}^2
 &=v_A^2\frac{\sin^2\alpha}{\cos^2\alpha}+v_A^2
      -2v_Av_N\frac{\sin^2\alpha}{\cos\alpha}\\
 &=\frac{v_A^2}{\cos^2\alpha}
      -2v_Av_N\frac{\sin^2\alpha}{\cos\alpha},
\end{align*}
\[
v_{RN}=\frac{v_A}{\cos\alpha}
\sqrt{1-2\frac{v_N}{v_A}\sin^2\alpha\cos\alpha}.
\]
The official final result is
\[
\boxed{\sin\beta=
\frac{\sin\alpha\cos\alpha}
{\sqrt{1-2(v_N/v_A)\sin^2\alpha\cos\alpha}}}.
\]
